\documentclass[12pt,a4paper]{ctexart}

\usepackage{geometry}
\geometry{left=2.5cm,right=2.5cm,top=2cm,bottom=2cm}
\usepackage{xcolor}
\usepackage{enumitem}
\usepackage{hyperref}

\setlength{\parindent}{0pt}
\setlength{\parskip}{0.75em}
\linespread{1.15}

\newcommand{\note}[1]{\textcolor{gray}{\small #1}}

\begin{document}

\begin{center}
{\Large \textbf{讲解讲稿 · 多种算法求解0/1背包问题的性能比较}}

\vspace{0.3cm}
讲解时长：约7--8分钟 \quad 幻灯片页数：16
\end{center}

\vspace{0.6cm}

\note{说明：本讲稿按新版PPT顺序撰写，语速正常时约7--8分钟。录制时不需要逐字照读，可以把每页加粗观点作为讲解重点。}

% ====== Slide 1 ======
\section*{Slide 1 — 封面}

各位老师、同学大家好。我们小组汇报的题目是“多种算法求解0/1背包问题的性能比较”。本实验属于2.3问题求解类，选取的是经典的0/1背包问题。我们主要实现并比较了四种算法：贪心算法、动态规划算法、回溯算法和分支限界算法，并在FSU、Pi和JJ三个数据源上进行了实验。接下来我会按照“问题定义—算法设计—数据集—实验结果—综合结论”的顺序进行讲解。

% ====== Slide 2 ======
\section*{Slide 2 — 0/1背包问题}

首先看0/1背包问题的定义。给定$n$个物品，每个物品有重量$w_i$和价值$v_i$，背包容量为$W$。每件物品只能选择一次，要么整件放入背包，要么不放入，不能像分数背包那样拆开。我们的目标是在不超过容量限制的前提下，使选中物品的总价值最大。

这个问题虽然形式很简单，但它是经典的NP-hard组合优化问题。也就是说，如果直接枚举所有可能的选择，一共有$2^n$种组合，规模稍微变大就会非常困难。同时，它又具有最优子结构，所以可以用动态规划等方法在一定条件下求出精确解。因此，它很适合用来比较不同算法策略在时间效率、空间消耗和解质量上的差异。

% ====== Slide 3 ======
\section*{Slide 3 — 四种求解策略}

这一页是四种算法的总体对比。贪心算法的思路最直接：按照价值重量比排序，优先选择当前看起来最划算的物品。它的时间复杂度主要来自排序，是$O(n\log n)$，但不保证最优。

动态规划算法通过状态转移逐步构造最优解，时间复杂度是$O(nW)$，可以保证最优，但对容量$W$非常敏感。

回溯算法把每个物品看作“选”或“不选”两个分支，深度优先搜索解空间树，并通过上界剪枝减少搜索。分支限界算法也是搜索解空间树，但它用优先队列优先扩展更有希望的节点。它们都可以得到最优解，但最坏情况下仍然是指数级复杂度。

所以这四种算法的瓶颈并不一样：贪心主要受数据分布影响，动态规划主要受容量影响，回溯和分支限界主要受物品数量影响。

% ====== Slide 4 ======
\section*{Slide 4 — 贪心算法}

贪心算法的具体做法是，先计算每个物品的单位重量价值，也就是$r_i=v_i/w_i$，然后按$r_i$从大到小排序。排序完成后，依次尝试把物品放进背包，只要当前剩余容量够就放入，不够就跳过。

它的优势非常明显：实现简单，运行速度快，并且几乎不受背包容量$W$影响。无论容量是一千、一百万还是更大，它主要做的仍然只是排序和一次线性扫描。

但是它的问题也很明显：贪心只关注当前局部最优，不能回退。PPT中的例子展示了这一点：如果先选择了一个性价比较高但占用空间较大的物品，可能会导致后面两个组合价值更高的物品无法同时放入。也就是说，贪心算法适合快速给出一个可行解，但不能保证全局最优。

% ====== Slide 5 ======
\section*{Slide 5 — 动态规划}

动态规划是本实验中最稳定的精确算法。我们定义$dp[w]$表示容量为$w$时能够获得的最大价值。对每个物品，状态转移为：
\[
dp[w]=\max(dp[w],dp[w-w_i]+v_i)
\]
这里使用一维数组进行空间优化，因此空间复杂度从二维形式降为$O(W)$。

实现时有一个关键细节：容量循环必须从大到小倒序遍历。因为0/1背包中每个物品只能使用一次，如果正序遍历，就可能在同一轮中重复使用同一件物品，问题会变成完全背包。

动态规划的优点是能够保证最优解，实验中它在所有可运行实例上都得到了最优值。但它的复杂度$O(nW)$是伪多项式复杂度，容量$W$越大，数组越大，更新次数也越多。例如Pi-L数据中，$n=2000$、容量约250万时，DP耗时已经达到10秒左右。因此，DP适合容量可控的精确求解场景。

% ====== Slide 6 ======
\section*{Slide 6 — 回溯与分支限界}

这一页是我负责的回溯和分支限界部分。回溯算法把0/1背包问题看成一棵二叉解空间树。对每个物品，我们有两个选择：选它或者不选它。深度优先遍历这棵树时，每一条从根到叶子的路径都代表一种物品组合。

如果完全枚举，最坏情况需要搜索$2^n$个节点，因此必须加入剪枝。本实验中使用了两类剪枝：第一类是可行性剪枝，如果当前重量已经超过容量，就停止扩展；第二类是上界剪枝，如果当前价值加上剩余物品的理论最大上界仍然不超过已知最优值，就直接剪掉这个分支。

分支限界算法和回溯的目标相同，都是寻找最优解，但搜索顺序不同。回溯更像“先沿着一条路走到底”，而分支限界用优先队列保存活节点，优先扩展上界最高、最有希望产生更优解的节点。这样通常能够更早找到较好的可行解，从而增强后续剪枝效果。不过它也有代价：优先队列可能保存大量节点，所以空间压力比普通回溯更大。

% ====== Slide 7 ======
\section*{Slide 7 — 三个数据源}

本实验使用了三个数据源。FSU数据集规模较小，包含8个实例，$n$从5到24，并且有官方最优解，因此主要用于验证算法正确性。

Pi数据集是我们按固定随机种子生成的实验数据，包括不相关、弱相关和强相关三种价值重量关系，每种关系又分为S、M、L三个规模。它的作用是控制变量，观察规模和数据分布对算法的影响。

JJ数据集选取了6个较大规模实例，$n$为400或1000，容量为100万。它更接近大规模测试场景，用于观察贪心、DP和搜索类算法在规模增大后的可运行性差异。

% ====== Slide 8 ======
\section*{Slide 8 — FSU：四种算法同表比较}

这一页展示FSU数据集中四种算法的求解值对比。可以看到，DP、回溯和分支限界在8个实例上都达到了最优值，说明三个精确算法的实现是正确的。贪心算法只有P01达到最优，其余实例都有不同程度的偏差。

其中最明显的是P06，贪心结果为1478，而最优值为1735，gap达到14.81\%。这个例子说明，即使物品数量只有7个，贪心也可能因为局部选择失误而错过全局最优组合。

所以FSU的意义不只是“数据小”，而是它很好地证明了：贪心算法的速度优势不能等价于求解质量优势，而DP、回溯和分支限界作为精确算法，在小规模基准上都能稳定得到最优解。

% ====== Slide 9 ======
\section*{Slide 9 — FSU：四种算法耗时比较}

这一页比较FSU上四种算法的运行时间。大多数小实例中，四种算法都在毫秒级甚至更低，差距并不大。但P08非常有代表性：它只有24个物品，但容量达到6,404,180。DP因为需要维护容量维度的状态数组，耗时达到757.697毫秒；而贪心只需要排序和扫描，耗时只有0.081毫秒；回溯和分支限界也因为$n$较小、剪枝有效，分别只用了0.306毫秒和1.378毫秒。

这个现象非常重要：DP的瓶颈是容量$W$，而回溯和分支限界的瓶颈是物品数$n$。所以在$n$小但$W$很大的场景下，搜索类算法反而可能比DP更快。这也解释了为什么不能简单地说某一种精确算法永远最好。

% ====== Slide 10 ======
\section*{Slide 10 — Pi-S：四种算法完整比较}

Pi-S是$n=50$的小规模生成数据，因此四种算法都可以完整运行。结果显示，DP、回溯和分支限界全部命中最优。贪心在不相关和弱相关S实例上也达到最优，但在强相关S实例上出现0.85\%的gap。

这里的关键在于数据分布。强相关数据中，物品价值和重量关系更接近，单位价值差别变小，贪心排序的区分度下降，所以更容易选错组合。反过来，不相关数据中单位价值差异更明显，贪心更容易识别高性价比物品。

从耗时看，搜索类算法在S规模下表现仍然可接受。比如强相关S实例中，回溯耗时0.218毫秒，分支限界耗时1.549毫秒，都能快速求得最优。但这只是小规模下的结果，不能直接推广到$n=500$或$n=2000$。

% ====== Slide 11 ======
\section*{Slide 11 — 搜索类算法：从可运行到不可运行}

这一页重点解释回溯和分支限界为什么在部分数据集中被标记为SKIP。在FSU和Pi-S中，物品数量较小，搜索树虽然理论上是指数级，但在剪枝后仍可以快速完成。

但是到了Pi-M，$n=500$；Pi-L，$n=2000$；以及JJ中$n=400$或1000，搜索空间会膨胀到无法完整运行的程度。即使有上界剪枝，也无法保证在合理时间内遍历足够多的节点。因此，我们没有把它们当作程序失败，而是在实验设计中主动记录为SKIP。

这反映了搜索类算法的真实适用边界：它们适合小规模精确求解，也适合展示剪枝思想；但面对几百到几千个物品时，指数级复杂度会迅速成为硬性限制。

% ====== Slide 12 ======
\section*{Slide 12 — Pi：规模扩展与贪心偏差}

这一页进一步看Pi数据的规模扩展。随着规模从S到M再到L，DP耗时明显上升。以L规模为例，三类数据的DP耗时都在10秒左右，其中强相关L达到11,416.532毫秒。这与$O(nW)$复杂度是一致的：物品数增加，容量也较大，状态更新次数自然增加。

贪心算法则始终保持毫秒级，最大gap为强相关S上的0.85\%。这说明在Pi生成数据中，贪心大多数时候能得到接近最优的结果，但它的风险仍然集中在性价比差异不明显的数据上。

这一页的结论是：DP更适合作为最优基准，贪心更适合作为快速近似方案；而回溯和分支限界只能在S规模参与完整比较，M和L规模不具备稳定可运行性。

% ====== Slide 13 ======
\section*{Slide 13 — JJ：大规模实例下的现实选择}

JJ数据集展示的是大规模实例下的现实选择。这里的$n$达到400或1000，容量为100万。DP仍然可以运行并获得最优解，但耗时大约在727到1632毫秒之间，明显高于小规模实例。贪心算法运行时间大多在毫秒级，gap最大只有0.04\%。

这说明，在JJ这类结构化大规模数据上，贪心虽然没有理论最优保证，但实际结果非常接近最优。另一方面，回溯和分支限界因为$n$太大，仍然被记录为SKIP。

因此，JJ实验给出的现实启示是：当数据规模较大时，如果必须保证最优，DP仍然是可选方案；如果更看重速度，贪心可以作为非常有效的近似方案；而搜索类精确算法在这种规模下不适合直接完整运行。

% ====== Slide 14 ======
\section*{Slide 14 — 核心洞察：$n$与$W$的博弈}

这一页是本实验的核心洞察。影响算法表现的并不是单一因素，而是$n$和$W$共同决定的。

当$n$小、$W$也小时，四种算法都能轻松运行。当$n$大但$W$较小时，DP更有优势，因为它的复杂度虽然和$n$有关，但只要$W$可控，就能稳定求最优。当$n$小但$W$很大时，回溯和分支限界可能优于DP，因为它们不按容量建表，而是主要受搜索树规模影响。当$n$和$W$都很大时，贪心往往是最现实的选择，因为它不受容量影响，且时间复杂度最低。

所以，实验最终不是为了证明某一个算法绝对最好，而是为了说明：不同算法有不同的“适用区域”。算法选型要看数据规模、容量大小和精度要求。

% ====== Slide 15 ======
\section*{Slide 15 — 四种算法复杂度对比}

这一页用表格总结四种算法的复杂度和适用性。贪心时间复杂度为$O(n\log n)$，空间复杂度$O(n)$，速度最快，但不保证最优，主要受数据分布影响。

DP时间复杂度为$O(nW)$，空间复杂度$O(W)$，能够保证最优，但容量越大，时间和内存压力越明显。

回溯最坏时间复杂度为$O(2^n)$，空间复杂度主要来自递归栈和路径记录，通常可看作$O(n)$级别；它保证最优，但只适合$n$较小的精确求解。分支限界同样最坏为$O(2^n)$，但通过优先队列和上界估计更有方向地搜索；它也保证最优，但最坏空间可能达到指数级。

综合来看，四种算法构成了一条从“快速近似”到“精确最优”的谱系。实际使用时，不能只看理论最优性，也要考虑能否在有限时间和内存中运行完成。

% ====== Slide 16 ======
\section*{Slide 16 — 总结}

最后总结一下本实验的主要结论。第一，0/1背包问题虽然定义简单，但非常适合展示不同算法策略之间的差异。第二，贪心算法最快，几乎不受容量影响，但不保证最优，在FSU P06中gap达到14.81\%。第三，动态规划结果最稳定，在所有可运行实例上都得到最优，但受容量$W$限制明显。第四，回溯和分支限界在小规模数据上也能得到最优解，并且在$n$小、$W$大的情况下可能比DP更快，但它们面对大规模$n$时会受到指数级搜索空间限制。第五，算法选择的关键不是问“哪个算法最好”，而是根据$n$、$W$和精度要求选择最合适的策略。

我们的最终理解是：时间效率与最优性的折中，是算法设计中非常核心的问题。以上就是我们的汇报，谢谢大家。

\end{document}
